\usepackage{fvextra}
\usepackage{colortbl}
\usepackage{booktabs}
\usepackage{etoolbox}
\usepackage{titlesec}
\usepackage{tcolorbox}
\tcbuselibrary{skins,breakable}
\usepgflibrary{svg.path}
\geometry{paperwidth=\PaperPageWidth mm,paperheight=\PaperPageHeight mm,left=\PaperPageMarginX mm,right=\PaperPageMarginX mm,top=\PaperPageMarginY mm,bottom=\PaperPageMarginY mm}
\newcommand{\PaperCurrentCodeScale}{\PaperCodeScale}
\newcommand{\PaperCodeFontSize}{\fontsize{\fpeval{\PaperBodySize*\PaperCurrentCodeScale}pt}{\fpeval{\PaperBodySize*\PaperCurrentCodeScale*\PaperBodyLeading}pt}\selectfont}
\fvset{fontsize=\PaperCodeFontSize,breaklines=true,breakanywhere=true,breaknonspaceingroup=true}
\newcommand{\PaperHeadingFont}[1]{\fontsize{\fpeval{\PaperBodySize*#1}pt}{\fpeval{\PaperBodySize*#1*\PaperHeadingLeading}pt}\selectfont}
\titleformat{\section}{\normalfont\bfseries\PaperHeadingFont{\PaperHeadingOne}}{\thesection}{1em}{}
\titleformat{\subsection}{\normalfont\bfseries\PaperHeadingFont{\PaperHeadingTwo}}{\thesubsection}{1em}{}
\titleformat{\subsubsection}{\normalfont\bfseries\PaperHeadingFont{\PaperHeadingThree}}{\thesubsubsection}{1em}{}
\titleformat{\paragraph}[block]{\normalfont\bfseries\PaperHeadingFont{\PaperHeadingFour}}{\theparagraph}{1em}{}
\titleformat{\subparagraph}[block]{\normalfont\bfseries\PaperHeadingFont{\PaperHeadingFive}}{\thesubparagraph}{1em}{}
\arrayrulecolor{paper-table-rule}
\setlength{\heavyrulewidth}{\PaperBorderWidth em}
\setlength{\lightrulewidth}{\PaperBorderWidth em}
\setlength{\tabcolsep}{\PaperTablePaddingX em}
\setlength{\extrarowheight}{\PaperTablePaddingY em}
\colorlet{shadecolor}{paper-code-bg}
\DefineVerbatimEnvironment{Highlighting}{Verbatim}{commandchars=\\\{\},breaklines=true,breakanywhere=true,fontsize=\PaperCodeFontSize}
\newlanguage\PaperNoHyphenation
\newcommand{\PaperInlineCodeBreaks}{\language\PaperNoHyphenation\let\PaperInlineOperatorTok\OperatorTok\renewcommand{\OperatorTok}[1]{\PaperInlineOperatorTok{##1}\allowbreak}}
\providecommand{\VERB}{}\renewcommand{\VERB}{\Verb[commandchars=\\\{\},breakanywhere=false,formatcom=\PaperInlineCodeBreaks]}
\renewcommand{\KeywordTok}[1]{\textcolor{paper-code-keyword}{\textbf{#1}}}
\renewcommand{\DataTypeTok}[1]{\textcolor{paper-code-type}{#1}}
\renewcommand{\ControlFlowTok}[1]{\textcolor{paper-code-keyword}{\textbf{#1}}}
\renewcommand{\FunctionTok}[1]{\textcolor{paper-code-name}{#1}}
\renewcommand{\VariableTok}[1]{\textcolor{paper-code-name}{#1}}
\renewcommand{\BuiltInTok}[1]{\textcolor{paper-code-name}{#1}}
\renewcommand{\AttributeTok}[1]{\textcolor{paper-code-name}{#1}}
\renewcommand{\ImportTok}[1]{\textcolor{paper-code-scope}{#1}}
\renewcommand{\ConstantTok}[1]{\textcolor{paper-code-name}{#1}}
\renewcommand{\StringTok}[1]{\textcolor{paper-code-string}{#1}}
\renewcommand{\CharTok}[1]{\textcolor{paper-code-string}{#1}}
\renewcommand{\SpecialCharTok}[1]{\textcolor{paper-code-string}{#1}}
\renewcommand{\SpecialStringTok}[1]{\textcolor{paper-code-string}{#1}}
\renewcommand{\VerbatimStringTok}[1]{\textcolor{paper-code-string}{#1}}
\renewcommand{\CommentTok}[1]{\textcolor{paper-code-comment}{\textit{#1}}}
\renewcommand{\CommentVarTok}[1]{\textcolor{paper-code-comment}{\textit{#1}}}
\renewcommand{\AlertTok}[1]{\textcolor{paper-code-error}{#1}}
\renewcommand{\ErrorTok}[1]{\textcolor{paper-code-error}{\textbf{#1}}}
\renewcommand{\WarningTok}[1]{\textcolor{paper-code-comment}{\textbf{\textit{#1}}}}
\renewcommand{\DocumentationTok}[1]{\textcolor{paper-code-comment}{\textit{#1}}}
\renewcommand{\InformationTok}[1]{\textcolor{paper-code-comment}{\textit{#1}}}
\renewcommand{\PreprocessorTok}[1]{\textcolor{paper-code-preprocessor}{#1}}
\renewcommand{\OperatorTok}[1]{\textcolor{paper-code-punctuation}{#1}}
\renewcommand{\ExtensionTok}[1]{\textcolor{paper-code-ellipsis}{\textbf{#1}}}
\renewcommand{\DecValTok}[1]{\textcolor{paper-code-number}{#1}}
\renewcommand{\BaseNTok}[1]{\textcolor{paper-code-number}{#1}}
\renewcommand{\FloatTok}[1]{\textcolor{paper-code-number}{#1}}
\newcommand{\PaperHighlight}[2]{\bgroup\markoverwith{\textcolor{#1}{\rule[-0.4ex]{2pt}{2.2ex}}}\ULon{#2}}
\newcommand{\PaperUnmarked}[1]{#1}
\newcommand{\PaperStrikeHighlight}[1]{\bgroup\markoverwith{\rlap{\textcolor{\PaperRemovalColor}{\rule[-0.4ex]{2pt}{2.2ex}}}\rule[0.55ex]{2pt}{0.4pt}}\ULon{#1}}
\begingroup\lccode`\~=`\ \lowercase{\endgroup\def\PaperCodeSpace{~}}
\ExplSyntaxOn
\cs_generate_variant:Nn \seq_set_split:Nnn {NVn}
\NewDocumentCommand{\PaperCodeStrike}{m}{\seq_set_split:NVn\l_tmpa_seq\PaperCodeSpace{#1}\seq_pop_left:NN\l_tmpa_seq\l_tmpa_tl\exp_args:NV\sout\l_tmpa_tl\seq_map_inline:Nn\l_tmpa_seq{\sout{\mbox{\space}}\allowbreak\sout{##1}}}
\NewDocumentCommand{\PaperDottedUrlText}{m}{\tl_map_inline:nn{#1}{\leavevmode\setbox0=\hbox{##1}\rlap{\raisebox{-\PaperLinkUnderlineOffset em}{\hbox to\wd0{\leaders\hbox{\rule{\PaperLinkUnderlineWidth em}{\PaperLinkUnderlineWidth em}\hskip\PaperLinkUnderlineWidth em}\hfill}}}\box0\penalty100\hskip0pt}}
\NewDocumentCommand{\PaperLinkText}{m}{\tl_if_in:nnTF{#1}{\CSLIndent}{\group_begin:\cs_set_eq:NN\PaperOriginalCSLIndent\CSLIndent\RenewDocumentCommand\CSLIndent{m}{\PaperOriginalCSLIndent{\PaperDottedUrlText{##1}}}#1\group_end:}{\PaperDottedLink{#1}}}
\ExplSyntaxOff
\newcommand{\PaperDiffLine}[2]{\expandafter\def\csname PaperDiffLine#1\endcsname{#2}}
\newcommand{\PaperDiffBackground}[1]{\ifcsname PaperDiffLine\arabic{FancyVerbLine}\endcsname\begingroup\setlength{\fboxsep}{0pt}\colorbox{\csname PaperDiffLine\arabic{FancyVerbLine}\endcsname}{#1}\endgroup\else#1\fi}
\newcommand{\PaperStrikeToken}[1]{\csletcs{PaperOriginal#1}{#1}\csdef{#1}##1{\csuse{PaperOriginal#1}{\PaperCodeStrike{##1}}}}
\newif\ifPaperDiff
\newif\ifPaperRemoved
\newenvironment{PaperComparison}{\begingroup\renewcommand{\PaperCurrentCodeScale}{\PaperComparisonCodeScale}}{\endgroup}
\newenvironment{PaperMeasureCell}[3]{\def\PaperMeasuredName{#1}\expandafter\newsavebox\csname PaperCell#1\endcsname\begin{lrbox}{\csname PaperCell#1\endcsname}\begin{minipage}[t]{\dimexpr#2\dimexpr\linewidth-#3\tabcolsep*2\relax\relax}\raggedright}{\end{minipage}\end{lrbox}\global\setbox\csname PaperCell\PaperMeasuredName\endcsname=\copy\csname PaperCell\PaperMeasuredName\endcsname}
\newcommand{\PaperCenteredCell}[3]{\begingroup\dimen0=\ht\csname PaperCell#2\endcsname\dimen2=\dp\csname PaperCell#2\endcsname\ifdim\ht\csname PaperCell#3\endcsname>\dimen0\dimen0=\ht\csname PaperCell#3\endcsname\fi\ifdim\dp\csname PaperCell#3\endcsname>\dimen2\dimen2=\dp\csname PaperCell#3\endcsname\fi\ifdim\ht\csname @arstrutbox\endcsname>\dimen0\dimen0=\ht\csname @arstrutbox\endcsname\fi\ifdim\dp\csname @arstrutbox\endcsname>\dimen2\dimen2=\dp\csname @arstrutbox\endcsname\fi\raisebox{\dimexpr(\dimen0-\dimen2-\ht\csname PaperCell#1\endcsname+\dp\csname PaperCell#1\endcsname)/2\relax}{\usebox{\csname PaperCell#1\endcsname}}\endgroup}
\newenvironment{PaperDiff}{\begingroup\PaperDifftrue}{\endgroup}
\let\PaperIconnote\PaperIconinfo
\let\PaperIcondraftnote\PaperIconinfo
\newcommand{\PaperCalloutIcon}[1]{%
  \resizebox{1.75em}{!}{\begin{tikzpicture}[yscale=-1]
    \path[use as bounding box] (-1pt,-1pt) rectangle (25pt,25pt);
    \pgfsetlinewidth{2pt}\pgfsetroundcap\pgfsetroundjoin
    \csname PaperIcon#1\endcsname\pgfusepath{stroke}
  \end{tikzpicture}}%
}
\newcommand{\PaperCalloutLabel}[1]{{\color{PaperCalloutAccent}\sffamily\small\bfseries\MakeUppercase{#1}}}
\newlength{\PaperPanelExtraInset}
\newenvironment{PaperNumberedGutter}{%
  \setlength{\PaperPanelExtraInset}{\PaperPanelInset em}%
}{}
\tcbset{paper panel/.style={
  enhanced, breakable, boxrule=0pt, frame hidden, arc=0.25em, sharp corners=west,
  left skip=\dimexpr\PaperPanelInset em+\PaperPanelExtraInset\relax, boxsep=0pt, left=\dimexpr\PaperPanelPadding em+\PaperPanelRule em\relax,
  right=\PaperPanelPadding em, top=\PaperPanelPadding em, bottom=\PaperPanelPadding em,
  before skip=0.7\baselineskip, after skip=0.7\baselineskip
}, paper gutter icon/.style 2 args={
  overlay unbroken and first={
    \node[anchor=north east,inner sep=0pt,text=#2]
      at ([xshift=\dimexpr-0.4em-\PaperPanelExtraInset\relax,yshift=-0.35em]frame.north west) {\PaperCalloutIcon{#1}};}
}}
\newtcolorbox{PaperCallout}[1]{
  paper panel, colback=paper-#1-bg, colupper=paper-text,
  borderline west={\PaperPanelRule em}{0pt}{paper-#1-accent}, paper gutter icon={#1}{paper-#1-accent},
  before upper={\colorlet{PaperCalloutAccent}{paper-#1-accent}%
    \setlength{\PaperPanelExtraInset}{0pt}%
    \setlength{\parskip}{0.6\baselineskip}%
    \renewenvironment{quote}{\begin{PaperCallout}{quote}}{\end{PaperCallout}}}
}
\newtcolorbox{PaperWording}{
  paper panel, colback=paper-wording-bg, colupper=paper-wording-text,
  borderline west={\PaperPanelRule em}{0pt}{paper-wording-border},
  paper gutter icon={wording}{paper-wording-border},
  before upper={\emergencystretch=\linewidth\setlength{\parskip}{0.6\baselineskip}%
    \setlength{\PaperPanelExtraInset}{0pt}}
}
\newenvironment{PaperAddition}{%
  \begin{PaperCallout}{insert}%
  \color{paper-add}\colorlet{shadecolor}{paper-insert-bg}\colorlet{paper-wording-add-bg}{paper-insert-bg}%
}{\end{PaperCallout}}
\newenvironment{PaperRemoval}{%
  \begin{PaperCallout}{remove}\PaperRemovedtrue\color{paper-rm}%
  \colorlet{shadecolor}{paper-remove-bg}\colorlet{paper-wording-remove-bg}{paper-remove-bg}%
}{\end{PaperCallout}}
\AtBeginEnvironment{Highlighting}{
  \ifPaperDiff
    \renewcommand{\NormalTok}[1]{\textcolor{paper-wording-text}{#1}}
    \renewcommand{\FancyVerbHighlightLine}[1]{\PaperDiffBackground{#1}}
    \renewcommand{\VariableTok}[1]{\textcolor{paper-add}{#1}}
    \renewcommand{\StringTok}[1]{\textcolor{paper-rm}{\PaperCodeStrike{#1}}}
  \fi
  \ifPaperRemoved
    \forcsvlist{\PaperStrikeToken}{NormalTok,KeywordTok,DataTypeTok,ControlFlowTok,FunctionTok,VariableTok,BuiltInTok,AttributeTok,ImportTok,ConstantTok,StringTok,CharTok,SpecialCharTok,SpecialStringTok,VerbatimStringTok,CommentTok,CommentVarTok,DocumentationTok,InformationTok,PreprocessorTok,OperatorTok,ExtensionTok,DecValTok,BaseNTok,FloatTok,AlertTok,ErrorTok,WarningTok}
  \fi
}
\newcommand{\PaperDottedLink}[1]{\bgroup\markoverwith{\rule[-\PaperLinkUnderlineOffset em]{\PaperLinkUnderlineWidth em}{\PaperLinkUnderlineWidth em}\hspace{\PaperLinkUnderlineWidth em}}\ULon{#1}}
\makeatletter
\AtBeginDocument{\let\PaperUrlInternal\url@\def\url@#1{\begingroup\def\Url@FormatString{\UrlFont\expandafter\PaperDottedUrlText\expandafter{\Url@String}}\PaperUrlInternal{#1}\endgroup}}
\AtBeginDocument{\NewCommandCopy\PaperHref\href\RenewDocumentCommand\href{o m m}{\IfNoValueTF{#1}{\PaperHref{#2}{\PaperLinkText{#3}}}{\PaperHref[#1]{#2}{\PaperLinkText{#3}}}}\NewCommandCopy\PaperHyperlink\hyperlink\RenewDocumentCommand\hyperlink{m m}{\PaperHyperlink{#1}{\PaperDottedLink{#2}}}\NewCommandCopy\PaperHyperref\hyperref\renewrobustcmd\hyperref{\@ifnextchar[{\PaperLocalRef}{\PaperHyperref}}\def\PaperLocalRef[#1]#2{\PaperHyperref[#1]{\PaperDottedLink{#2}}}}
\AtBeginDocument{\patchcmd{\@maketitle}{\LARGE}{\PaperHeadingFont{\PaperHeadingOne}}{}{}}
\AtBeginDocument{\pretocmd{\@makecaption}{\begingroup\color{paper-caption}}{}{}\apptocmd{\@makecaption}{\endgroup}{}{}}
\AtBeginDocument{\@ifundefined{LT@makecaption}{}{\patchcmd{\LT@makecaption}{\reset@font}{\reset@font\color{paper-caption}}{}{}}}
\renewcommand{\pnum}[1]{\noindent\makebox[0pt][l]{\makebox[0pt][r]{\textcolor{paper-wording-number}{\scriptsize #1}\hspace{\dimexpr\@totalleftmargin+1.4em\relax}}}\ignorespaces}
\AtBeginDocument{
  \fontsize{\PaperBodySize pt}{\fpeval{\PaperBodySize*\PaperBodyLeading}pt}\selectfont
  \color{paper-text}
  \global\let\default@color\current@color
  \hypersetup{linkcolor=paper-link,urlcolor=paper-link,citecolor=paper-link,pdfborder={0 0 0},pdfborderstyle={},linkbordercolor=paper-link,urlbordercolor=paper-link,citebordercolor=paper-link}
}
\makeatother
